\chapter*{คำนำ}
\addcontentsline{toc}{chapter}{คำนำ}

คู่มือวิชาการและชุดพรอมพต์การเรียนรู้เล่มนี้ จัดทำขึ้นเพื่อเป็นแนวทางเชิงปฏิบัติการสำหรับนักศึกษา นักวิจัย เกษตรกรรุ่นใหม่ และผู้สนใจเริ่มต้นศึกษาการพัฒนานวัตกรรมเกษตรดิจิทัล โดยมุ่งเน้นการบูรณาการองค์ความรู้ข้ามศาสตร์ 4 เสาหลัก ได้แก่ ทฤษฎีฟิสิกส์เชิงหลักการแรก ระบบสมองกลฝังตัวและการสื่อสารอุตสาหกรรม การเรียนรู้เชิงลึกระดับอุปกรณ์ และการพัฒนาแอปพลิเคชันสมาร์ตโฟนด้วยเทคโนโลยีสมัยใหม่

ในสภาพแปลงเกษตรกรรมจริง การตรวจวัดสมบัติของดินมักได้รับผลกระทบจากความผันผวนของสภาพแวดล้อม โดยเฉพาะการเลื่อนลอยตามอุณหภูมิและความไม่เป็นเชิงเส้นของหัววัด คู่มือเล่มนี้จึงได้นำเสนอระเบียบวิธี \textbf{Physics-Informed Neural Network (PINN)} ซึ่งฝังกฎเกณฑ์ทางฟิสิกส์ เช่น สภาพยอมทางไฟฟ้าและการปรับเทียบตามทฤษฎีอาร์เรเนียส ลงในโครงข่ายประสาทเทียม เพื่อให้สมาร์ตโฟนสามารถชดเชยค่าความคลาดเคลื่อนได้อย่างแม่นยำในระดับเสี้ยววินาที

คณะผู้จัดทำหวังเป็นอย่างยิ่งว่า คู่มือและชุดพรอมพต์วิศวกรรมที่ได้รับการจัดทำอย่างเป็นระบบนี้ จะเป็นประโยชน์ต่อการจัดการเรียนการสอน การวิจัย และการสร้างสรรค์นวัตกรรมด้านเกษตรแม่นยำของประเทศให้ก้าวหน้าอย่างยั่งยืนสืบไป

\vspace{1.5cm}
\begin{flushright}
    ผู้ช่วยศาสตราจารย์ ดร.ชีวะ ทัศนา\\
    สาขาวิชาฟิสิกส์ คณะวิทยาศาสตร์และเทคโนโลยี\\
    มหาวิทยาลัยราชภัฏรำไพพรรณี
\end{flushright}
